\documentclass[10pt]{article}
\usepackage[T1]{fontenc}
\usepackage[utf8]{inputenc}
\usepackage[french]{babel}
\usepackage{amsmath,amssymb,mathrsfs,mathpazo}
\usepackage{geometry}
\geometry{a4paper,margin=2cm}
\usepackage{graphicx,array,tabularx,tabularray,colortbl,multirow,multicol,booktabs,enumitem,pifont,ulem,fancybox,diagbox,cancel}
\usepackage{pstricks,pstricks-add,pst-plot,pst-tree,pst-node,pst-text,pst-func,pst-math,pst-bspline,pst-3dplot,pst-eucl,pst-solides3d}
\let\pscalculate\undefined
\usepackage{pst-bezier}
\usepackage{tikz,pgfplots,tkz-tab}
\usetikzlibrary{arrows,arrows.meta,patterns,calc,positioning,trees,decorations.markings,shapes.geometric,shapes.misc,shapes.arrows,intersections,decorations.pathmorphing,decorations.pathreplacing,backgrounds,shadows,babel}
\pgfplotsset{compat=1.18}
\usepgfplotslibrary{fillbetween}
\usepackage[np]{numprint}
\usepackage{listings,makecell,pas-tableur,twoopt}
\let\barre\undefined
\usepackage[pstricks]{tablvar}
\definecolor{darkgreen}{RGB}{0,100,0}
\providecommand{\red}{\color{red}}
\providecommand{\blue}{\color{blue}}
\providecommand{\green}{\color{green}}
\providecommand{\black}{\color{black}}
\providecommand{\white}{\color{white}}
\providecommand{\gray}{\color{gray}}
\providecommand{\cyan}{\color{cyan}}
\providecommand{\magenta}{\color{magenta}}
\providecommand{\yellow}{\color{yellow}}
\providecommand{\violet}{\color{violet}}
\providecommand{\orange}{\color{orange}}
\providecommand{\euro}{\texteuro}
\providecommand{\EUR}{\texteuro}
\providecommand{\ieme}{\textsuperscript{e}}
\providecommand{\degres}{\ensuremath{^\circ}}
\providecommand{\degree}{\ensuremath{^\circ}}
\providecommand{\vv}[1]{\overrightarrow{#1}}
\providecommand{\Coord}[1]{#1}
\providecommand{\decosix}{}
\providecommand{\starredbullet}{\ensuremath{\bullet}}
\pagestyle{empty}
\setlength{\parindent}{0pt}
\usepackage{adjustbox}\DeclareUnicodeCharacter{2236}{:}\DeclareUnicodeCharacter{2190}{\ensuremath{\leftarrow}}
\providecommand{\R}{\mathbb{R}}
\providecommand{\N}{\mathbb{N}}
\providecommand{\Z}{\mathbb{Z}}
\providecommand{\Q}{\mathbb{Q}}
\providecommand{\D}{\mathbb{D}}
\providecommand{\C}{\mathbb{C}}
\providecommand{\e}{{\mathrm{e}}}
\providecommand{\vect}[1]{\overrightarrow{#1}}
\providecommand{\vectt}[1]{\overrightarrow{\mathrm{#1}}}
\providecommand{\vvt}[1]{\overrightarrow{\mathrm{#1}}}
\providecommand{\np}[2][]{{#2~#1}}
\providecommand{\vv}[1]{{\overrightarrow{#1}}}
\renewcommand{\d}{{\mathrm{d}}}
\providecommand{\degres}{{^\circ}}
\providecommand{\degre}{{^\circ}}
\providecommand{\blue}{\color{blue}}
\providecommand{\red}{\color{red}}
\providecommand{\enskip}{\quad}
\providecommand{\llbracket}{\lbrack\!\lbrack}
\providecommand{\rrbracket}{\rbrack\!\rbrack}
\providecommand{\euro}{{\text{€}}}
\providecommand{\no}{{\text{n°}}}
\providecommand{\interval}[2]{{[#1;#2]}}
\providecommand{\intervFO}[2]{{[#1;#2[}}
\providecommand{\dotfill}{\ldots}

\providecommand{\ds}{\displaystyle}
\providecommand{\textding}[1]{#1}
\providecommand{\Oij}{\(\mathrm{O};\vec{\imath},\vec{\jmath}\)}
\providecommand{\Oijk}{\(\mathrm{O};\vec{\imath},\vec{\jmath},\vec{k}\)}
\providecommand{\Ouv}{\(\mathrm{O};\vec{u},\vec{v}\)}
\providecommand{\intervFF}[2]{\ensuremath{[#1;#2]}}
\renewcommand{\euro}{\ifmmode\text{\texteuro}\else\texteuro\fi}
\DeclareUnicodeCharacter{00B0}{\ifmmode^{\circ}\else\textdegree\fi}
\DeclareUnicodeCharacter{20AC}{\ifmmode\text{\texteuro}\else\texteuro\fi}
% Énoncés : Éducation nationale. Transcription et corrigés : APMEP, Denis Vergès et contributeurs. Découpage et classement : CoopMaths, Sébastien Lozano, Éric Elter et contributeurs. Adaptation de la présentation : Maths Explorer.
% CC BY-SA : https://coopmaths.fr/www/about/
% https://coopmaths.fr/alea/static/e3c/2020/tex/e3c_2020_00_sujet2_4_cor.tex
\begin{document}
\noindent\textit{APMEP / CoopMaths. CC BY-SA. Présentation adaptée par Maths Explorer.}\par\medskip
\section*{Exercice 4 (5 points)}
	

\begin{center}
	\psset{unit=0.75cm,arrowsize=2pt 4}
	\begin{pspicture}(17.5,5.8)
		\psframe[fillstyle=solid,fillcolor=lightgray](17.5,5.8)
		\psline[linecolor=white,linewidth=1.5pt](0,1)(14.3,5.8)
		\psline[linecolor=white](1.35,1.4)(2.5,0.8)(9,3)(7.7,3.5)
		\psline(3.3,2.1)(3.3,3.2)(5.85,4.1)(5.85,2.9)(9.7,4.2)(16.7,1)(5.85,2.9)
		\psline(16.7,1)(3.3,2.1)
		\psline[linecolor=white,linestyle=dashed](3.3,2.1)(2.5,2.5)(5,3.4)(5.85,2.9)
		\psline[linestyle=dashed](3.3,3.2)(2.5,3.7)(5,4.5)(5.85,4.1)
		\psline[linestyle=dashed](2.5,2.5)(2.5,3.7)
		\psline[linestyle=dashed](5,3.4)(5,4.5)
		\psline{<->}(5.6,3)(9.6,4.4)\uput[u](7.6,3.7){9 m}
		\psline{<->}(10,4.3)(16.7,1.2)\uput[u](13.35,2.75){18 m}
		\psline(9.45,4.07)(9.7,3.95)(10,4.1)
		\uput[d](3.3,2){A}\uput[d](5.85,2.9){B}\uput[u](9.7,4.2){D}\uput[dr](16.7,1){T}
	\end{pspicture}
\end{center}
	
	\subsection*{1.}
	Les droites $(TD)$ et $(AB)$ sont perpendiculaires, les vecteurs $\overrightarrow{TD}$ et $\overrightarrow{DB}$ sont donc orthogonaux : leur produit scalaire est donc nul.
	
	\subsection*{2.}
	D’après la relation de Chasles : $\overrightarrow{TA} \cdot \overrightarrow{TB} = (\overrightarrow{TD} + \overrightarrow{DA}) \cdot \overrightarrow{TB} = \overrightarrow{TD} \cdot \overrightarrow{TB} + \overrightarrow{DA} \cdot \overrightarrow{TB}$.
	
	Or $\overrightarrow{TD} \cdot \overrightarrow{TB} = \overrightarrow{TD} \cdot \overrightarrow{TD} = TD^2 = 324$.
	
	D’autre part $\overrightarrow{DA} \cdot \overrightarrow{TB} = \overrightarrow{DA} \cdot \overrightarrow{DB} = (7,32 + 9) \times 9 = 146,88$.
	
	Finalement :
	\[
	\overrightarrow{TA} \cdot \overrightarrow{TB} = 324 + 146,88 = 470,88
	\]
	
	\subsection*{3. Méthode 1:}
	On sait que $\overrightarrow{TA} \cdot \overrightarrow{TB} = TA \times TB \times \cos (\widehat{ATB}) \quad(1)$.
	
	Or d’après le théorème de Pythagore appliqué au triangle $TDB$ rectangle en $D$, on a :
	\[
	AT^2 = AD^2 + TD^2 = 16,322 + 182 = 590,3424.
	\]
	L’égalité $(1)$ s’écrit donc :
	\[
	470,88 = \sqrt{405} \times \sqrt{590,3424} \times \cos (\widehat{ATB}), \text{ d’où }
	\cos (\widehat{ATB}) = \frac{470,88}{\sqrt{405} \times \sqrt{590,3424}} \approx 0,96301.
	\]
	
	La calculatrice donne $\widehat{ATB} \approx 15,63^{\circ}, \text{ soit } \widehat{ATB} \approx 15,6^{\circ} \text{ au dixième près.}$
	
	\subsection*{3. Méthode 2:}
	Dans le triangle $ATD$ rectangle en $D$ :\\
	$ \tan (\widehat{ATD}) = \frac{AD}{TD} = \frac{9+7,32}{18} \approx 0,90667, \text{ d’où } \widehat{ATD} \approx 42,1975°.$\\
	Dans le triangle $BTD$ rectangle en $D$ :\\
	$ \tan (\widehat{BTD}) = \frac{BD}{TD} = \frac{9}{18} = \frac{1}{2}, \text{ d’où } \widehat{BTD} \approx 26,5651^{\circ}.$\\
	On a donc par différence $\widehat{ATB} \approx 42,1975^{\circ} - 26,5651^{\circ} \approx 15,6324^{\circ} \text{ soit } 15,6^{\circ} \text{ au dixième près.}$

\end{document}
